\chapter{FPGA Implementation}\label{chap:fpga-impl}
This chapter details the VHDL design and Vivado integration.

\section{Architecture Overview}
The design comprises parameterized `multiplier` units, `neuron` blocks aggregating weighted sums with bias, and a top-level `nn\_rgb` module orchestrating layers. The `sigmoid\_lut`/`sigmoid\_IP` provides activation using a precomputed lookup table stored in block RAM.

\section{VHDL Modules}
Key files in `VHDL/` include:
\begin{itemize}
    \item `config.vhd`: Global configuration (bit widths, scaling factors).
    \item `multiplier.vhd`: Fixed-point multiply stage using DSP resources.
    \item `neuron.vhd`: Accumulation and bias addition; interfaces to activation.
    \item `nn\_rgb.vhd`: Top entity connecting layers and streaming interfaces.
    \item `sigmoid\_lut.vhd` and `sigmoid\_IP.vhd`: Activation lookup.
    \item `tb\_nn\_rgb.vhd`, `sim\_nn\_rgb.vhd`: Testbench and simulation harness.
\end{itemize}

\section{Simulation}
ModelSim/Questa simulation artifacts are under `sim/`. Provided test vectors allow functional verification against golden outputs. Waveforms and logs (e.g., `vsim.wlf`) support debugging.

\section{Synthesis and Implementation}
The Vivado project `vivado_thesis/thesis_ultrascale_1/` targets an UltraScale device. Post-implementation reports (timing, utilization, clocking) are stored under `thesis_ultrascale_1.runs/`. The design is clocked to meet real-time throughput for streaming pixels.

\section{Resource/Performance Considerations}
Design parameters (bit widths, layer sizes) affect LUTs, FFs, BRAM, and DSP usage. Pipelining and parallelization trade latency for throughput. Activation via LUT reduces DSP demand versus CORDIC or piecewise-linear implementations.
